% !TEX root = ../main.tex
\section{Modelamiento del problema}

El problema consiste en seleccionar, para cada elemento constructivo optimizable
de una vivienda, exactamente una alternativa desde el catálogo técnico
correspondiente a su categoría. Cada alternativa posee una transmitancia
térmica y un costo por metro cuadrado. El objetivo es minimizar el costo total
de los materiales seleccionados, manteniendo la transmitancia térmica promedio
ponderada por área de la vivienda por debajo de un valor máximo permitido
$U_{\max}$.

Debido a que la decisión asociada a cada alternativa es del tipo
seleccionar/no seleccionar, el problema se formula como un modelo de
Programación Lineal Entera Binaria.

\subsection{Conjuntos e índices}

Se definen los siguientes conjuntos:

\begin{itemize}
    \item $C$: conjunto de categorías de elementos constructivos de la vivienda,
    por ejemplo, muros perimetrales, techumbre, pisos, ventanas y puertas.
    \item $I$: conjunto de elementos constructivos de la vivienda.
    \item $I^{O} \subseteq I$: conjunto de elementos optimizables o libres.
    \item $I^{F} \subseteq I$: conjunto de elementos fijos, cuya solución
    constructiva ya fue definida previamente por el usuario.
    \item $J_c$: conjunto de alternativas disponibles en el catálogo para la
    categoría $c \in C$.
\end{itemize}

Además, se utiliza la función $g(i)$ para indicar la categoría a la cual
pertenece cada elemento $i \in I$. Por lo tanto, las alternativas admisibles
para un elemento $i$ son las contenidas en $J_{g(i)}$.

Los índices utilizados son:

\begin{itemize}
    \item $i$: índice de elementos constructivos.
    \item $c$: índice de categorías.
    \item $j$: índice de alternativas de material o solución constructiva.
\end{itemize}

\subsection{Parámetros}

Los parámetros utilizados por el modelo se presentan en la
Tabla~\ref{tab:parametros}.

\begin{table}[h]
\centering
\caption{Parámetros del modelo}
\label{tab:parametros}
\begin{tabular}{|p{2.2cm}|p{5.3cm}|p{2.3cm}|p{3.2cm}|}
\hline
\textbf{Parámetro} & \textbf{Descripción} & \textbf{Unidad} & \textbf{Fuente / rango} \\ \hline
$A_i$ & Área del elemento constructivo $i$. & m$^2$ & Generador de instancias. Debe ser positiva. \\ \hline
$U_{ij}$ & Transmitancia térmica de la alternativa $j$ aplicable al elemento optimizable $i$. & W/m$^2$K & Catálogo técnico generado para la instancia. \\ \hline
$C_{ij}$ & Costo por metro cuadrado de la alternativa $j$ aplicable al elemento optimizable $i$. & u.m./m$^2$ & Catálogo técnico generado para la instancia. \\ \hline
$U_i^{F}$ & Transmitancia térmica del material predefinido para el elemento fijo $i$. & W/m$^2$K & Dato de la instancia. \\ \hline
$C_i^{F}$ & Costo por metro cuadrado del material predefinido para el elemento fijo $i$. & u.m./m$^2$ & Dato de la instancia. \\ \hline
$U_{\max}$ & Transmitancia térmica promedio ponderada máxima permitida para la vivienda. & W/m$^2$K & Parámetro configurable de cada instancia. \\ \hline
$A_T$ & Área total considerada en el balance térmico, definida como $A_T=\sum_{i\in I}A_i$. & m$^2$ & Calculada a partir de las áreas de la instancia. \\ \hline
\end{tabular}
\end{table}

El enunciado no fija rangos numéricos específicos para las áreas, costos,
transmitancias ni para $U_{\max}$. Estos valores deben ser definidos por el
generador de instancias de forma técnicamente coherente. En particular, para
las alternativas de una misma categoría debe respetarse la relación
cualitativa indicada en el proyecto: alternativas con menor transmitancia
térmica, es decir, con mejor aislación, deben presentar un mayor costo por
metro cuadrado.

Respecto al tamaño de las instancias, el proyecto establece:

\begin{itemize}
    \item Instancias pequeñas: 3--4 categorías, 5--10 alternativas por
    categoría y 2--4 elementos libres.
    \item Instancias medianas: 5 categorías, 11--25 alternativas por
    categoría y 5--8 elementos libres.
    \item Instancias grandes: 5 categorías, 26--50 alternativas por
    categoría y 9--15 elementos libres.
\end{itemize}

Las dimensiones generales de la vivienda, tales como superficie habitable,
altura y volumen interior, forman parte de la caracterización de cada
instancia. En el modelo propuesto, su efecto se incorpora a través de las
áreas $A_i$ asignadas a cada elemento constructivo.

\subsection{Variable de decisión}

Para cada elemento optimizable $i \in I^{O}$ y cada alternativa
$j \in J_{g(i)}$, se define la variable binaria:

\[
x_{ij}=
\begin{cases}
1, & \text{si se selecciona la alternativa } j \text{ para el elemento } i,\\
0, & \text{en otro caso.}
\end{cases}
\]

Por lo tanto,

\[
x_{ij}\in\{0,1\},
\qquad
\forall i\in I^{O},\;
\forall j\in J_{g(i)}.
\]

No se define una variable de decisión para los elementos pertenecientes a
$I^{F}$, debido a que su material se encuentra fijado de antemano y no puede
ser modificado por el modelo.

\subsection{Función objetivo}

El objetivo es minimizar el costo total de materiales de la vivienda. Para los
elementos optimizables, el costo de una alternativa corresponde al costo por
metro cuadrado multiplicado por el área del elemento. Los elementos fijos
aportan un costo constante.

La función objetivo queda definida como:

\[
\min Z =
\sum_{i\in I^{O}}
\sum_{j\in J_{g(i)}}
A_i C_{ij}x_{ij}
+
\sum_{i\in I^{F}}
A_i C_i^{F}.
\]

El primer término representa el costo de las soluciones constructivas
seleccionadas por el modelo para los elementos optimizables. El segundo
término corresponde al costo de los elementos fijos.

Como el segundo término es constante, no modifica cuál combinación de
alternativas resulta óptima. Sin embargo, se mantiene en la formulación para
que $Z$ represente el costo total de materiales de la vivienda. Si en la
implementación los elementos fijos se consideran una inversión ya realizada
o costo hundido, dicho término puede omitirse sin cambiar la solución
óptima de los elementos libres.

\subsection{Restricciones}

\subsubsection{Selección de una alternativa por elemento optimizable}

Cada elemento libre debe utilizar exactamente una alternativa del catálogo
correspondiente a su categoría:

\[
\sum_{j\in J_{g(i)}} x_{ij}=1,
\qquad
\forall i\in I^{O}.
\]

Esta restricción evita que un elemento optimizable quede sin material
asignado y también impide que se seleccionen simultáneamente dos o más
alternativas para el mismo elemento.

\subsubsection{Desempeño térmico global}

La transmitancia térmica promedio ponderada por área de la vivienda debe
ser menor o igual que el límite $U_{\max}$.

Considerando tanto elementos optimizables como fijos, la transmitancia promedio
puede expresarse como:

\[
\overline{U} =
\frac{
\displaystyle
\sum_{i\in I^{O}}
\sum_{j\in J_{g(i)}}
A_iU_{ij}x_{ij}
+
\sum_{i\in I^{F}}
A_iU_i^{F}
}{
\displaystyle
\sum_{i\in I}A_i
}.
\]

Por requerimiento del problema:

\[
\overline{U}\leq U_{\max}.
\]

Dado que el denominador corresponde al área total de la vivienda y es un
parámetro conocido, la restricción anterior puede escribirse de manera
lineal como:

\[
\sum_{i\in I^{O}}
\sum_{j\in J_{g(i)}}
A_iU_{ij}x_{ij}
+
\sum_{i\in I^{F}}
A_iU_i^{F}
\leq
U_{\max}
\sum_{i\in I}A_i.
\]

De esta forma, se conserva la estructura lineal del modelo y se garantiza que
el desempeño térmico global de la vivienda no exceda el umbral permitido.

\subsection{Modelo matemático completo}

El modelo puede resumirse de la siguiente manera:

\[
\boxed{
\begin{aligned}
\min \quad
& Z=
\sum_{i\in I^{O}}
\sum_{j\in J_{g(i)}}
A_iC_{ij}x_{ij}
+
\sum_{i\in I^{F}}
A_iC_i^{F}
\\[2mm]
\text{s.a.}\quad
& \sum_{j\in J_{g(i)}}x_{ij}=1,
&& \forall i\in I^{O},
\\[2mm]
&
\sum_{i\in I^{O}}
\sum_{j\in J_{g(i)}}
A_iU_{ij}x_{ij}
+
\sum_{i\in I^{F}}
A_iU_i^{F}
\leq
U_{\max}
\sum_{i\in I}A_i,
\\[2mm]
& x_{ij}\in\{0,1\},
&& \forall i\in I^{O},
\quad \forall j\in J_{g(i)}.
\end{aligned}
}
\]

\subsection{Supuestos del modelamiento}

Para representar el problema mediante Programación Lineal Entera Binaria se
adoptan los siguientes supuestos:

\begin{itemize}
    \item Cada elemento constructivo pertenece a una única categoría.
    \item Cada elemento optimizable debe recibir exactamente una alternativa
    disponible para su categoría.
    \item Las alternativas del catálogo son compatibles con la categoría
    para la cual fueron generadas.
    \item El área de cada elemento es conocida y permanece constante durante
    la optimización.
    \item La transmitancia térmica y el costo por metro cuadrado de cada
    alternativa son datos conocidos de la instancia y no dependen de las
    demás decisiones.
    \item Los costos son proporcionales al área cubierta por cada material.
    \item El aporte térmico global se obtiene mediante una suma ponderada por
    área, por lo que se asume aditividad entre los elementos.
    \item Los elementos fijos no pueden ser modificados por el modelo, pero su
    transmitancia se considera dentro del balance térmico global.
    \item El valor de $U_{\max}$ es conocido antes de resolver cada instancia.
    \item No se consideran en esta entrega costos adicionales de instalación,
    mano de obra, transporte, disponibilidad de stock ni interacciones
    constructivas entre materiales, ya que el enunciado no entrega información
    para modelarlos.
    \item Las instancias entregadas al modelo deben ser factibles, es decir,
    debe existir al menos una selección de alternativas que satisfaga todas
    las restricciones.
\end{itemize}

La formulación cumple los supuestos de proporcionalidad y aditividad propios
de los modelos lineales. Además, la naturaleza binaria de $x_{ij}$ permite
representar directamente la selección discreta de materiales.
